\IfFileExists{stacks-project.cls}{%
\documentclass{stacks-project}
}{%
\documentclass{amsart}
}
\usepackage{amssymb}

% For dealing with references we use the comment environment
\usepackage{verbatim}
\newenvironment{reference}{\comment}{\endcomment}
%\newenvironment{reference}{}{}
\newenvironment{slogan}{\comment}{\endcomment}
\newenvironment{history}{\comment}{\endcomment}

% For commutative diagrams we use Xy-pic
\usepackage[all]{xy}

% We use 2cell for 2-commutative diagrams.
\xyoption{2cell}
\UseAllTwocells

% We use multicol for the list of chapters between chapters
\usepackage{multicol}

% This is generally recommended for better output
\usepackage{lmodern}
\usepackage[T1]{fontenc}
\usepackage{textalpha}

% For cross-file-references
\usepackage{xr-hyper}
\newcommand{\maybeexternaldocument}[2]{%
  \IfFileExists{#2.aux}{%
    \externaldocument[#1]{#2}%
  }{%
    \PackageWarning{stack}{Missing #2.aux; external references with prefix #1 unavailable}%
  }%
}
\let\externaldocumentorig\externaldocument
\renewcommand{\externaldocument}[2][]{%
  \IfFileExists{#2.aux}{%
    \externaldocumentorig[#1]{#2}%
  }{%
    \PackageWarning{stack}{Missing #2.aux; external references with prefix #1 unavailable}%
  }%
}

% Package for hypertext links:
\usepackage{hyperref}

% For any local file, say "hello.tex" you want to link to please
% use \externaldocument[hello-]{hello}
\externaldocument[introduction-]{introduction}
\externaldocument[conventions-]{conventions}
\externaldocument[sets-]{sets}
\externaldocument[categories-]{categories}
\externaldocument[topology-]{topology}
\externaldocument[sheaves-]{sheaves}
\externaldocument[sites-]{sites}
\externaldocument[stacks-]{stacks}
\externaldocument[fields-]{fields}
\externaldocument[algebra-]{algebra}
\externaldocument[brauer-]{brauer}
\externaldocument[homology-]{homology}
\externaldocument[derived-]{derived}
\externaldocument[simplicial-]{simplicial}
\externaldocument[more-algebra-]{more-algebra}
\externaldocument[smoothing-]{smoothing}
\externaldocument[modules-]{modules}
\externaldocument[sites-modules-]{sites-modules}
\externaldocument[injectives-]{injectives}
\externaldocument[cohomology-]{cohomology}
\externaldocument[sites-cohomology-]{sites-cohomology}
\externaldocument[dga-]{dga}
\externaldocument[dpa-]{dpa}
\externaldocument[sdga-]{sdga}
\externaldocument[hypercovering-]{hypercovering}
\externaldocument[schemes-]{schemes}
\externaldocument[constructions-]{constructions}
\externaldocument[properties-]{properties}
\externaldocument[morphisms-]{morphisms}
\externaldocument[coherent-]{coherent}
\externaldocument[divisors-]{divisors}
\externaldocument[limits-]{limits}
\externaldocument[varieties-]{varieties}
\externaldocument[topologies-]{topologies}
\externaldocument[descent-]{descent}
\externaldocument[perfect-]{perfect}
\externaldocument[more-morphisms-]{more-morphisms}
\externaldocument[flat-]{flat}
\externaldocument[groupoids-]{groupoids}
\externaldocument[more-groupoids-]{more-groupoids}
\externaldocument[etale-]{etale}
\externaldocument[chow-]{chow}
\externaldocument[intersection-]{intersection}
\externaldocument[pic-]{pic}
\externaldocument[weil-]{weil}
\externaldocument[adequate-]{adequate}
\externaldocument[dualizing-]{dualizing}
\externaldocument[duality-]{duality}
\externaldocument[discriminant-]{discriminant}
\externaldocument[derham-]{derham}
\externaldocument[local-cohomology-]{local-cohomology}
\externaldocument[algebraization-]{algebraization}
\externaldocument[curves-]{curves}
\externaldocument[resolve-]{resolve}
\externaldocument[models-]{models}
\externaldocument[functors-]{functors}
\externaldocument[equiv-]{equiv}
\externaldocument[pione-]{pione}
\externaldocument[etale-cohomology-]{etale-cohomology}
\externaldocument[proetale-]{proetale}
\externaldocument[relative-cycles-]{relative-cycles}
\externaldocument[more-etale-]{more-etale}
\externaldocument[trace-]{trace}
\externaldocument[crystalline-]{crystalline}
\externaldocument[spaces-]{spaces}
\externaldocument[spaces-properties-]{spaces-properties}
\externaldocument[spaces-morphisms-]{spaces-morphisms}
\externaldocument[decent-spaces-]{decent-spaces}
\externaldocument[spaces-cohomology-]{spaces-cohomology}
\externaldocument[spaces-limits-]{spaces-limits}
\externaldocument[spaces-divisors-]{spaces-divisors}
\externaldocument[spaces-over-fields-]{spaces-over-fields}
\externaldocument[spaces-topologies-]{spaces-topologies}
\externaldocument[spaces-descent-]{spaces-descent}
\externaldocument[spaces-perfect-]{spaces-perfect}
\externaldocument[spaces-more-morphisms-]{spaces-more-morphisms}
\externaldocument[spaces-flat-]{spaces-flat}
\externaldocument[spaces-groupoids-]{spaces-groupoids}
\externaldocument[spaces-more-groupoids-]{spaces-more-groupoids}
\externaldocument[bootstrap-]{bootstrap}
\externaldocument[spaces-pushouts-]{spaces-pushouts}
\externaldocument[spaces-chow-]{spaces-chow}
\externaldocument[groupoids-quotients-]{groupoids-quotients}
\externaldocument[spaces-more-cohomology-]{spaces-more-cohomology}
\externaldocument[spaces-simplicial-]{spaces-simplicial}
\externaldocument[spaces-duality-]{spaces-duality}
\externaldocument[formal-spaces-]{formal-spaces}
\externaldocument[restricted-]{restricted}
\externaldocument[spaces-resolve-]{spaces-resolve}
\externaldocument[formal-defos-]{formal-defos}
\externaldocument[defos-]{defos}
\externaldocument[cotangent-]{cotangent}
\externaldocument[examples-defos-]{examples-defos}
\externaldocument[algebraic-]{algebraic}
\externaldocument[examples-stacks-]{examples-stacks}
\externaldocument[stacks-sheaves-]{stacks-sheaves}
\externaldocument[criteria-]{criteria}
\externaldocument[artin-]{artin}
\externaldocument[quot-]{quot}
\externaldocument[stacks-properties-]{stacks-properties}
\externaldocument[stacks-morphisms-]{stacks-morphisms}
\externaldocument[stacks-limits-]{stacks-limits}
\externaldocument[stacks-cohomology-]{stacks-cohomology}
\externaldocument[stacks-perfect-]{stacks-perfect}
\externaldocument[stacks-introduction-]{stacks-introduction}
\externaldocument[stacks-more-morphisms-]{stacks-more-morphisms}
\externaldocument[stacks-geometry-]{stacks-geometry}
\externaldocument[moduli-]{moduli}
\externaldocument[moduli-curves-]{moduli-curves}
\externaldocument[examples-]{examples}
\externaldocument[exercises-]{exercises}
\externaldocument[guide-]{guide}
\externaldocument[desirables-]{desirables}
\externaldocument[coding-]{coding}
\externaldocument[obsolete-]{obsolete}
\externaldocument[fdl-]{fdl}
\externaldocument[index-]{index}

% Heap Project documents
\externaldocument[linear-algebra-]{linear-algebra}
\externaldocument[groups-]{groups}
\externaldocument[complex-analysis-]{complex-analysis}
\externaldocument[functional-analysis-]{functional-analysis}
\externaldocument[categories-]{categories}
\externaldocument[commutative-algebra-]{commutative-algebra}
\externaldocument[ringed-spaces-]{ringed-spaces}
\externaldocument[homological-algebra-]{homological-algebra}
\externaldocument[lie-algebras-]{lie-algebras}
\externaldocument[representation-theory-]{representation-theory}
\externaldocument[smooth-manifolds-]{smooth-manifolds}
\externaldocument[differential-topology-]{differential-topology}
\externaldocument[algebraic-topology-]{algebraic-topology}
\externaldocument[differential-geometry-]{differential-geometry}
\externaldocument[algebraic-geometry-]{algebraic-geometry}
\externaldocument[symplectic-geometry-]{symplectic-geometry}
\externaldocument[complex-geometry-]{complex-geometry}
\externaldocument[hodge-theory-]{hodge-theory}
\maybeexternaldocument{deformations-}{deformations}
\externaldocument[vertex-algebras-]{vertex-algebras}
\externaldocument[springer-]{springer}
\externaldocument[infty-categories-]{infty-categories}
\externaldocument[seminars-differential-geometry-]{%
seminars-differential-geometry}
\externaldocument[seminars-complex-geometry-]{%
seminars-complex-geometry}
\externaldocument[seminars-algebraic-geometry-]{%
seminars-algebraic-geometry}
\externaldocument[seminars-others-]{%
seminars-others}
%\externaldocument[geometric-prequantization-]{geometric-prequantization}
\externaldocument[surfaces-]{surfaces}
\externaldocument[birational-maps-commutative-algebra-]{birational-maps-commutative-algebra}
\externaldocument[cremona-transformations-]{cremona-transformations}
\externaldocument[derived-categories-of-sheaves-]{derived-categories-of-sheaves}
\externaldocument[geometric-invariant-theory-]{geometric-invariant-theory}
\externaldocument[geometric-stability-conditions-and-group-actions-]{geometric-stability-conditions-and-group-actions}
\externaldocument[moduli-spaces-of-sheaves-]{moduli-spaces-of-sheaves}
\externaldocument[bridgeland-stability-general-theory-]{bridgeland-stability-general-theory}
\externaldocument[hilbert-schemes-]{hilbert-schemes}
\externaldocument[noncommutative-abelian-surfaces-and-kummer-type-hyperkahler-varieties-]{noncommutative-abelian-surfaces-and-kummer-type-hyperkahler-varieties}
\externaldocument[mumford-tate-groups-in-hodge-theory-]{mumford-tate-groups-in-hodge-theory}
\externaldocument[hodge-birational-atoms-2026-]{hodge-birational-atoms-2026}
\externaldocument[xxii-egd-2026-]{%
xxii-egd-2026}
\externaldocument[pde-]{pde}
\externaldocument[probability-]{probability}

\externaldocument[linguistics-]{linguistics}
\externaldocument[genealogia-familiar-]{%
genealogia-familiar}

\externaldocument[%
16th-alga-meeting-2026-]{%
16th-alga-meeting-2026}

\externaldocument[%
iv-jornadas-lefschetz-2026-]{%
iv-jornadas-lefschetz-2026}

\externaldocument[reading-the-masters-]{%
reading-the-masters}

\externaldocument[real-analysis-]{%
real-analysis}

\externaldocument[optimal-transport-]{%
optimal-transport}

\externaldocument[history-]{history}

% Theorem environments.
%
\theoremstyle{plain}
\newtheorem{theorem}[subsection]{Theorem}
\newtheorem{proposition}[subsection]{Proposition}
\newtheorem{lemma}[subsection]{Lemma}

\theoremstyle{definition}
\newtheorem{definition}[subsection]{Definition}
\newtheorem{example}[subsection]{Example}
\newtheorem{exercise}[subsection]{Exercise}
\newtheorem{situation}[subsection]{Situation}

\theoremstyle{remark}
\newtheorem{remark}[subsection]{Remark}
\newtheorem{remarks}[subsection]{Remarks}

\numberwithin{equation}{subsection}

% Macros
%
\def\lim{\mathop{\mathrm{lim}}\nolimits}
\def\colim{\mathop{\mathrm{colim}}\nolimits}
\def\Spec{\mathop{\mathrm{Spec}}}
\def\Hom{\mathop{\mathrm{Hom}}\nolimits}
\def\Ext{\mathop{\mathrm{Ext}}\nolimits}
\def\SheafHom{\mathop{\mathcal{H}\!\mathit{om}}\nolimits}
\def\SheafExt{\mathop{\mathcal{E}\!\mathit{xt}}\nolimits}
\def\Sch{\mathit{Sch}}
\def\Mor{\mathop{\mathrm{Mor}}\nolimits}
\def\Ob{\mathop{\mathrm{Ob}}\nolimits}
\def\Sh{\mathop{\mathit{Sh}}\nolimits}
\def\NL{\mathop{N\!L}\nolimits}
\def\CH{\mathop{\mathrm{CH}}\nolimits}
\def\proetale{{pro\text{-}\acute{e}tale}}
\def\etale{{\acute{e}tale}}
\def\QCoh{\mathit{QCoh}}
\def\Ker{\mathop{\mathrm{Ker}}}
\def\Im{\mathop{\mathrm{Im}}}
\def\Coker{\mathop{\mathrm{Coker}}}
\def\Coim{\mathop{\mathrm{Coim}}}

% Boxtimes
%
\DeclareMathSymbol{\boxtimes}{\mathbin}{AMSa}{"02}

%
% Macros for moduli stacks/spaces
%
\def\QCohstack{\mathcal{QC}\!\mathit{oh}}
\def\Cohstack{\mathcal{C}\!\mathit{oh}}
\def\Spacesstack{\mathcal{S}\!\mathit{paces}}
\def\Quotfunctor{\mathrm{Quot}}
\def\Hilbfunctor{\mathrm{Hilb}}
\def\Curvesstack{\mathcal{C}\!\mathit{urves}}
\def\Polarizedstack{\mathcal{P}\!\mathit{olarized}}
\def\Complexesstack{\mathcal{C}\!\mathit{omplexes}}
% \Pic is the operator that assigns to X its picard group, usage \Pic(X)
% \Picardstack_{X/B} denotes the Picard stack of X over B
% \Picardfunctor_{X/B} denotes the Picard functor of X over B
\def\Pic{\mathop{\mathrm{Pic}}\nolimits}
\def\Picardstack{\mathcal{P}\!\mathit{ic}}
\def\Picardfunctor{\mathrm{Pic}}
\def\Deformationcategory{\mathcal{D}\!\mathit{ef}}


\begin{document}

\title{History}
\maketitle

\phantomsection
\label{section-phantom}

\tableofcontents

\section{Napoleón en España y en Nueva España
la retirada de Hidalgo}
\label{section-retirada-de-hidalgo}

% The original title, byline, and body below
% are transcribed verbatim from the DOCX.
% Only LaTeX markup and whitespace differ.
% The attribution footnote is editorial.

L. Barjau\footnote{Luis Barjau.
Texto reproducido íntegramente del documento
\emph{La retirada de Hidalgo.docx}.}

En 1807 el ejército de Napoleón Bonaparte
ingresó a España buscando la alianza para
invadir Portugal, quien se negó a cumplir
bloqueo comercial contra Gran Bretaña.
Napoleón quería someter a los ingleses
mediante guerra financiera.

Al año siguiente Fernando VII el rey Borbón
ascendió al trono de España. Pero fue
expulsado por Napoleón, que nombró a su
hermano José I al reinado que duró cinco
años, hasta 1813. Fernando y su padre, Carlos
IV, tuvieron que abdicar y fueron hechos
prisioneros en Francia. Este hecho socavó los
cimientos de la monarquía española lo cual
repercutió en todo el mundo hispano. Las
Juntas regionales y descentralizadas se
negaron a reconocer la legitimidad de los
Bonaparte y este movimiento se extendió hasta
las colonias americanas. En 1809 hubo
rebeliones en los Andes. Se establecieron
tribunales generales tanto en España como en
América. En la península, la Junta Central
fue empujada hacia el sur, en 1809 se
estableció en Sevilla, pero de allí tuvo que
huir a Cádiz donde recibió la protección de
naves británicas. Se establecieron Cortes que
incluyeron los reinos de América y con ellas
España obtuvo mucho mayor control del
habitual. Los españoles radicados en las
colonias rechazaron el concepto de igualdad
que las cortes consideraban y los criollos
prefirieron que los \emph{indios}
permanecieran controlados por España. Este
conflicto también fue un antecedente de la
Independencia. Los españoles de América se
mantuvieron en la actitud de que indígenas y
afrodescendientes no estuvieran representados
en las cortes españolas. Las cortes eran los
órganos legislativos que representaban a los
sectores de todos los reinos ibéricos frente
al Rey. Institución nacida desde la Edad
Media.

La actitud de los españoles americanos fue un
gesto revelador de la voluntad europea
respecto del Otro, el distinto, del
“salvaje”, del extraño a sus costumbres. Que
habría de permear toda la relación política y
cultural de Europa con América. Toda la
mentalidad del hombre americano. Del criollo
enriquecido, con su resentimiento contra el
peninsular español; de las castas, el pueblo,
los pobres, la confusión identitaria.

La posible representación había estado
liderada en Quito por José Mejía Lequerica
que propuso igualdad para indios y mestizos
excluyendo a africanos.

Hasta 1812 el término “americano”, refería
solo a los blancos europeos. A partir de ese
año se aceptó tal designación también para
afrodescendientes, indígenas y mestizos, lo
que así quedó consignado en la Constitución
de Cádiz, la primera de España. Y esta
influyó aunque de manera ambigua en la
emancipación mexicana. Pues introdujo los
principios democráticos que sembraron las
ideas de la soberanía y ello estimuló las
tendencias independentistas. Surgieron los
ayuntamientos constitucionales y las
elecciones locales. Por primera vez existía
el voto y la autogestión que combatiría al
absolutismo real, de haberse ejercido al pie
de la letra. En 1820, restablecida la
Constitución en España, el clero, militares y
terratenientes novohispanos temieron la
pérdida de sus privilegios por esas nuevas
leyes de las Cortes. En la guerra de
independencia mexicana algunos personajes
realistas pactaron con Vicente Guerrero con
la intención de mantener el poder. Que
aceptaban las leyes de Cádiz pero excluían
aquellas de orden decididamente liberales.
Diputados novohispanos como Miguel Ramos
Arizpe viajaron a Cádiz y su experiencia
impulsó la Constitución mexicana de 1824.

En 1810 se formaron Juntas en Caracas, Buenos
Aires y Santiago. Hubo un vacío de poder
tanto en España como en las colonias, que
propició la creación de tales Juntas, que
tomaban el control de las diferentes
regiones, porque la resistencia en España al
mandato de Bonaparte estaba fracasando. En
Sudamérica Mariano Moreno asumió el liderazgo
que intentó una Confederación que se
inspiraba en la revolución de independencia
de EUA. Los virreinatos hispanoamericanos
trataron de evitar estas Juntas. El 20 de
julio de 1810 se promulgó la Declaración de
Independencia de Colombia. Pero desde el 1 de
febrero de 1804, gracias a los líderes
Toussaint Louverture y Jean Jacques
Dessalines se había proclamado la
independencia de Haití. Y aún antes, en 1776
Estados Unidos se había independizado. El
devenir de las paradojas: España y Francia
fueron los primeros aliados de este
movimiento de liberación en Norteamérica.

El 16 de septiembre de 1810, con el grito de
Dolores, Miguel Hidalgo y Costilla dio inicio
a la revolución de Independencia en Nueva
España, la cual, después de once años de
lucha se declararía su triunfo el 27 de
septiembre de 1821.

Una larga batalla de diez años, como el
periodo antiguo de la guerra de Troya. Sólo
que en América la guerra estuvo animada
violentamente también por el rencor indígena
porque Occidente había truncado el desarrollo
de la milenaria cultura mesoamericana. Las
batallas fueron activadas por la turbamulta
desarmada, enajenada a una rabia histórica
que prefería estrellarse contra el ejército
novohispano hasta el sacrificio y la muerte.
Retumbaba en la guerra la voluntad de
eliminar a los infiltrados españoles,
despojarlos, matarlos. Miguel Hidalgo no
calculó al Polifemo que había despertado de
su antiguo sueño en las cuevas de
Chicomoztoc.

A medida que se dirigía la revolución hacia
el centro de la Nueva España el pelotón se
engrosaba como bola de nieve en picada.
Contingente de 80,000 campesinos con armas
improvisadas. Destrozar las nuevas ciudades
hispanas, liquidar a los peninsulares que
mandaban a diestra y siniestra y se
enriquecían con los bienes naturales y con la
explotación de la mano de obra.

Hidalgo no cesó de asombrarse ante la
proporción del fenómeno que había despertado.
Y cuando llegó a las puertas de la ciudad, en
La Marquesa, Monte de las Cruces (no obstante
después de haber sido recibido en Toluca con
la misma pompa de recepción que se daba a los
virreyes) y ya al borde de la gran ciudad
capital, habiendo ganado la batalla, ¡tomó la
retirada!

Abrumadora superioridad numérica la de los
independentistas, contra el ejército
virreinal organizado y al mando de Torcuato
Trujillo y Rodríguez, que había peleado en
Alcalá la Real, de Jaén, contra Napoleón
Bonaparte. Aunque el ejército realista fuera
derrotado, el choque desigual dejó un saldo
campesino de cerca de 4000 muertos.

Sin explicar nada a nadie el cura Hidalgo
regresó por el camino que lo había conducido
hasta el centro del poder hispano. “Viva
Fernando VII”, había entremezclado un año
antes en su grito inicial por la
Independencia. “Viva nuestra señora de
Guadalupe, viva Fernando VII, muera el mal
gobierno” fue la proclama completa. Confusión
y desconcierto. A favor del rey español, a
favor del cristianismo sincretizado en la
guadalupana; pero contra del poder virreinal.
¿Qué se buscaba? ¿Qué era posible? El nombre
del rey todavía implorado contra la invasión
francesa en España.

La complejidad de cualquier proceso
histórico, en la Nueva España concitó su
extremada y densa problemática. Si la
invasión napoleónica en la Península hispana
estimuló la posibilidad de que en sus reinos
americanos se observara la opción
independiente; que los criollos vislumbraran
un nacionalismo inicial; que los indígenas
exacerbaran una caótica intervención
vengativa, caótica porque a la vez afirmaban
el derecho a un catolicismo dogmático
amparado en el mito sincrético de la virgen
guadalupana; que el cura Hidalgo encabezara
el surgimiento de un fervor revolucionario
que se gestaba desde la ideología pietista
cristiana, a la vez que en la observación de
la posibilidad de la riqueza local, pero sin
la enorme extracción de bienes y servicios
ejercida por la corona española, la tragedia
mexicana de la guerra independentista habría
de durar 11 años.

La decisión de retirada en el Monte de las
Cruces la había de pagar muy cara el párroco
de Dolores. En Acatita de Baján, Coahuila, lo
tomarían prisionero; fusilado el 30 de julio
de 1811 en el patio del antiguo hospital Real
de Loreto, hoy palacio de gobierno de
Chihuahua. Otra historia hubiera sido si
Hidalgo avanzara hasta la capital. Para bien
o para mal.

Diez largos años de guerra y todas las
traiciones y fusilamientos de
independentistas, que conocemos. Fantasmas y
leyendas. Iturbide y “el abrazo de
Acatempan”. Iturbide era el Jacob de México.

Con la independencia, el país se revolvió con
su desesperado intento de convertirse en una
nación. ¿Volver al ámbito espiritual de la
cultura indígena? ¿Construir cimientos
nacionales, con cuáles modos de producción?
¿Erguirse en la explotación del subsuelo? ¿La
agricultura? ¿El comercio? ¿Como las naciones
europeas?

En las raíces de México tuvieron que arraigar
las severas contradicciones. En extremo
complejas. Sin que nadie tuviera plena
conciencia de ello. ¿Acaso Juárez fuera el
primero en vislumbrar un camino cierto? El
indígena de Guelatao que subió hasta la silla
presidencial el 21 de enero de 1858; 37 años
después de haberse consumado la
independencia. ¡Juárez! ¿Qué pensó
íntimamente? ¡Mucho daríamos por haber
conocido dicho pensamiento! Por haber podido
discutir en aquellos términos. Juárez íntimo.
Juárez secreto. El milagro de la nación
mexicana. ¿La extensión del efluvio
prehispánico? La esencia revolucionaria de
esta nación: cambiar, cambiar a toda costa,
buscar un camino de esplendor para el viejo
dolor mexicano que se arrastra por la rivera
del rio de su historia.

\end{document}
